\documentclass[../../main.tex]{subfiles}


\begin{document}
\chapter*{前言}
\setlength{\epigraphwidth}{0.8\textwidth}
\epigraph{Ludwig Boltzman, who spent much of his life studying statistical mechanics, died in 1906, by his own hand. Paul Ehrenfest, carrying on the work, died similarly in 1933. Now it is our turn to study statistical mechanics. Perhaps it will be wise to approach the subject cautiously.}{D.L. Goodstein}
\vspace{1.5em}

大约一年前，2025年的11月低，距离热统期末考试还有一个多月，我开始认真地自己学习统计物理。如玻尔兹曼和他的学生埃伦菲斯特，我立马感受到了巨大的痛苦。其中原因主要是教材选择的错误，因为我只是理论的学习者，我不需要像理论的创建者一样收到人们的诘难\footnote{玻尔兹曼一生都在为原子论而奔波，这样一个如今看来显而易见的事情在当时完全不被接受，直到二十世纪初玻尔兹曼去世后人们才逐渐通过无可辩驳的实验接受原子论}，在更换教材后，我的痛苦减轻不少。但是统计物理中有很多晦涩的地方，对于一个初学者，任何教材都显得有些过于难懂。在进过一个月的不间断琢磨和推导后，我自认为对统计物理有了不错的掌握，同时我也萌生了写下我的想法的想法，我想写一点关于统计物理的东西，那些正式的教科书不屑于或不方便写的内容，以帮助其他学习统计物理的人减少初学时的痛苦。

于是我便开始着手实行这一想法，但统计物理又开始肘击我了。我立马发现我对统计物理实际上根本不了解，在一年前一个月的学习中，我的重点放在了各个公式本身的推导上，但对于各物理概念的认识非常匮乏\footnote{由于考试的压力}。于是我又开始学习统计物理，在断断续续的学习后，终于到了2026年的暑假，我留在学校正式开始写。起初我本想连带写写热力学，但我很快发现我对很多热力学概念没有什么把握，写出来无非也是抄书。于是我放弃这个额外的想法，只写统计物理。2026年8月12日，我正是开始写我对统计物理的各种想法，此后一直到2026年11月18日，我几乎把所有的精力都放在了写想法这个事上。

很难说我写的到底是一个什么属性的东西，称其为笔记但它并不是对某一个教材或课堂的记录，称其为讲义又显得过于繁琐不够精简。但是我想暂且还是称其为讲义吧。因为我的初衷在于记录自己的想法和减轻其他同学的学习痛苦，所以这本讲义的一个特点就是格外细致。对每一个统计原理的阐述、每一个统计模型的提出、每一个公式的推导都非常详细。可以说这绝不是一本好的教材，因为它太过于细致导致读者失去了思考的机会。以公式为例，几乎所有的公式读者都可以看着讲义口算，几乎就是把草稿纸上的推导步骤照搬到了讲义上。另外过于细致的阐述本质上只是我本人的想法，读者接收后会很大程度上减少自己的想法，这是很不好的\footnote{我本人并不排斥推导简略的教材，只要逻辑清晰，哪怕引入非常生硬我也能接受。因为任何理论的发展都是以年为单位的，试图在几天学明白一个理论同时追求所谓的自然是不可能的}。可以说，学习时的痛苦是学习的必经之路。这份讲义的目的是记录我的想法和减轻别人的学习痛苦，它并不是一本值得反复阅读的好教材，甚至它不一定能保证读者学好统计物理。

本讲义还有一个特点是有不少的例题，它们穿插在正文中，对整个讲义有很大的作用。讲义中的例题并不是传统的针对某个理论的检验性习题，而是正文的有必要补充，只是这种补充不便于以正文形式存在，读者有必要认真阅读每一个例题。

我花了不少心思写下这份讲义，但其中的缺陷并不少。比如对某些热力学量的计算特别是熵，本讲义有时候会采取比较繁琐的所谓通法，它确实是通法，但在有些时候它并不方便，读者在正文自行体会。关于顺磁性物质的统计理论，本讲义的推导出现了两个错误。讲义中从\(\dd U=T\dd S-VM\dd B+\mu\dd N\)出发得到\(\dd F=-S\dd T-VM\dd B+\mu\dd N\)，又根据\(F=-\k T\ln Z\)得到磁化强度的统计表达式。错误之处在于，磁性系统正确的热力学基本微分方程为\(\dd U=T\dd S+VB\dd M+\mu\dd N\)，这给出特性函数\footnote{\(\tilde{F}\)被称为磁自由能，自由能仅对热力学基本微分方程中的\((T,S)\)做勒让德变换，磁自由能同时对\((T,S)\)和\((M,B)\)两对变量做勒让德变换}\(\dd \tilde{F}=-S\dd T-VM\dd B+\mu\dd N\)，而配分函数对应的正确热力学量是\(\tilde{F}=-\k T\ln Z\)，这样两个错误刚好抵消，我们得到了正确的磁化强度统计表达式。处了这两处较为明显的缺陷外，讲义里还缺少对不少内容的讨论。除去缺少的一大块相互作用系统和统计物理的前沿应用，单说讲义的重点无相互作用系统就有不少没讨论的细节。最后本讲义没有课后习题，这是一个很大的缺陷，缺少了引导读者思考的部分。

但不论怎么说，这份讲义汇聚了我对统计物理大部分的想法，写这份讲义的两个初衷我已达到。后续我还会再次修改、补充已有的内容，并补上热力学和相互作用系统的统计物理。现在我要稍稍放下统计物理，去完成我的本科课业任务了。

在前言部分，我想有必要向读者披露我在写作过程中关于AI的使用情况。首先，关于统计物理的学习，我使用了不少AI，一般都是我苦思不得的问题询问AI，有些得到了解决，有些并没有。但请读者放心，正文的所有文字都是我用键盘一个字一个字敲出来的，内容本身都产自我的大脑，哪怕是由AI告诉我的内容，我也找到了资料去考证。其次正文的部分图片，是由AI编写Python代码或者\tikzname 代码画出来的，我对画图并不在行。最后，笔者看不到的部分即讲义的\LaTeX 排版，这部分AI深度参与，特别是一些排版细节，但这些都不是内容本身。请读者放心，本讲义传达给读者的内容，完全由我本人产生或者由我从现存资料中学习而来。

对于阅读本讲义的读者，请自行决定阅读的多少以及其中内容的接受程度。整本讲义只是我的个人见解，读者可以据此讲义备考，但考试不会以此来出题，其中内容上的缺失或多余跟我本人没有任何关系，一切都取决于读者本人。当然，倘若读者对讲义有任何想法，欢迎与我邮箱\footnote{QQ邮箱：2944052508@qq.com}交流。

\vspace{1.5em}
\hfill
\begin{tabular}{c}
李祯\\
前卫南区\\
二〇二六年十一月
\end{tabular}

\end{document}